This appendix collects technical background relevant to the proposed spectral
construction. It is intended as a reference for readers less familiar with
operator theory and spectral analysis.

\subsection{Basic Operator Theory}

Let $\hs$ be a complex Hilbert space with inner product $\langle \cdot, \cdot \rangle$.

\begin{itemize}
  \item A \emph{linear operator} $T \colon \mathcal{D}(T) \to \hs$ is a linear
        map defined on a dense subspace $\mathcal{D}(T) \subset \hs$.
  \item $T$ is \emph{self-adjoint} if $\mathcal{D}(T) = \mathcal{D}(T^*)$ and
        $T = T^*$.
  \item The \emph{spectrum} of $T$, denoted $\sigma(T)$, is the set of $\lambda \in \C$
        for which $T - \lambda I$ is not invertible (as a bounded operator from
        $\hs$ to $\hs$).
  \item A self-adjoint operator has real spectrum: $\sigma(T) \subset \R$.
\end{itemize}

\begin{remark}
  Self-adjointness is a strictly stronger condition than symmetry ($T \subset T^*$).
  Many operator constructions in spectral approaches to RH fail at the step of
  establishing self-adjointness.
\end{remark}

\subsection{Spectral Theorem}

For a self-adjoint operator $T$ on $\hs$, the spectral theorem guarantees the
existence of a projection-valued measure $E$ such that
\[
  T = \int_\R \lambda \, dE(\lambda).
\]
Trace-class operators have a well-defined trace $\mathrm{tr}(T) = \int_\R \lambda \, d\mu_T(\lambda)$
where $\mu_T$ is the spectral measure.

\subsection{Trace Formulas}

The Selberg trace formula and related trace formulas express geometric or spectral
data as a sum over eigenvalues. In approaches to RH via spectral theory, the goal
is to construct an operator whose trace formula reproduces the explicit formula
for $\zeta$ or its logarithmic derivative.

The explicit formula for $\zeta$ relates sums over primes to sums over zeros:
\[
  \sum_p \log p \cdot p^{-s} = -\frac{\zeta'(s)}{\zeta(s)} - (\text{explicit terms}).
\]
A spectral construction matching this would give the zeros as spectral data.

\todo{The relationship between the proposed operator $H$ and any explicit formula
should be worked out explicitly when $H$ is fully defined.}

\subsection{Functional Analysis Prerequisites}

The following concepts may be needed in the construction of $H$:

\begin{itemize}
  \item Hardy spaces $H^2(\C^+)$ on the upper half-plane.
  \item Paley--Wiener theory relating Fourier transform support to analytic continuation.
  \item The theory of Krein--Feller operators or Schrödinger operators on intervals.
  \item Mellin transforms and their relationship to $\zeta$.
\end{itemize}

\todo{As the construction progresses, the specific functional-analytic tools needed
will be identified and their prerequisites will be stated here.}
